\chapter{反抗焰}
《魔圆》系列电影中有一个反复出现的意象，我们此前略微提到过：
草坡上两把白椅并排而立，小圆与焰并肩坐着。
在前两部剧场版的片头字幕中，她们相拥依偎，画面甜美、可爱而无邪；
而到了第三部，这一意象很快便失却了甜味。

《叛逆》第二篇的高潮中，焰在化为魔女的过程中重返这一场景。
然而这一次，小圆起身站在椅子上，侧身倒下，在草地上溅作一片粉红色的痕迹。
而焰只能徒劳地伸出手，试图抓住她。
焰蹲在小圆身旁，因惊骇而圆睁双眼，四周则围着一群高挑而瘦长的“焰”，低头俯视着她。
随即，一个狂怒的“焰”挥出巨拳，将这个蹲在小圆残迹旁哭喊的焰击碎。% TODOfig 01:22:13

小圆已经不在，原本连贯统一的自我身份被一种弥散的抽象所取代。
焰失败了。
此刻，焰在审判焰，并判定她自己不够格。
积压已久的愤怒与悲恸终于倾泻而出，她粉碎了自己的身份，也把自己变成一种抽象而幽秘的存在：
魔女。

如同沙耶香，也很可能如同此前的每一个魔女，焰的魔女形态是一轮无尽的自我鞭笞。
这是一出心理剧：她不断重演将自己逼入绝望的事件，再为自己的失败施加惩罚。
她举枪自戕，而中弹的那个“自己”化作了她曾出于怜悯、不得不亲手杀死的小圆。
她不能死，也不配像小圆那样死去，因为她没能救下小圆。

问题还不止于此：焰正是小圆消失的原因。
焰一次次跨越时间的循环，使小圆获得成为圆环之理的力量，也正因此，小圆被从现实中抹除。
她与丘比谈论小圆一事，给了孵化者们足够的信息来构建如今正朝其收紧的圈套。
它们又正是借由焰来完成这一圈套。
焰是小圆最大的负累。

焰的魔女形态是最为直白的一种。
尖顶黑帽、突出的鼻子和下巴，除去数百英尺高的骷髅身躯之外，她几乎就是标准的万圣节女巫造型。
焰明知魔女为何物、从何而来，却依然未能躲开圈套，甚至为挫败丘比而主动投入其中。
沙耶香自视为骑士，因此即使化为魔女，仍保留着骑士的形貌；
而焰却清楚自己选择成为的是什么。
同样地，她也是有意牺牲自己。
正如她对丘比所言，她相信麻美和杏子会杀死自己。
于是，她的使魔将她引向断头台：那既是牺牲的装置，也是审判她罪行的刑具。

同时，围绕焰的还有一整套胡桃夹子意象。
一类使魔是长着胡桃夹子般颌部的巨牙\footnote{%
    译注：Luiselotte 的职责是驱逐老鼠。
    她们是狩猎白色老鼠的蛀牙骑兵，其中一部分也协助锡兵队驱逐无理阻碍送葬之人。
}；
另一类使魔形似玩具士兵\footnote{%
    译注：Lotte 的职责是执行处决。
    她们组成送葬队伍，将魔女带往断头台；憎恨愚者的愚者，也会对其无情断罪。
    笨重的头脑不知变通，完全不听魔女所言；它们最讨厌白色老鼠。
}，高耸的毛皮帽又使它们近似传统圣诞胡桃夹子。
霍夫曼《胡桃夹子与老鼠王》的原作中，孩子们一眼便认出胡桃夹子是士兵\footnote{%
    本章与下一章凡涉及《胡桃夹子与老鼠王》故事之处，均据 L. R. C. 的英译本\cite{hoffmann-nodate-nutcracker-mouse-king}。
}。
焰刚刚开始意识到自己正是将魔法少女困在迷宫中的魔女时，画面中出现了一张咧嘴而笑、用牙齿夹住核桃的嘴。% TODOfig
她的头颅失去了一半，只余下颌，这也映照着霍夫曼故事及柴可夫斯基\footnote{%
    译注：彼得·伊里奇·柴可夫斯基（1840年5月7日—1893年11月6日），俄国作曲家。
    其许多作品在有生之年就享誉国际，至今仍是最重要的浪漫主义音乐代表之一。
}据之创作的著名芭蕾舞剧中的同名胡桃夹子，后者失去的正是下颌。
许多使魔的玩偶质感，尤其是补充资料中名为“克拉拉”的布偶\footnote{%
    译注：Clara Dolls 的职责是哭泣。
    她们是来自伪街的孩子们，假扮成泪水的换装人偶，令葬礼更加喧闹。
    到场的人偶包括：自大、阴沉、欺骗、冷血、自私、毒舌、愚钝、嫉妒、懒惰、虚荣、懦弱、愚蠢、自卑与固执；
    最后的“爱”却未曾到来。
    她们拥有不亚于魔法少女的力量。
}\cite{magicaquartet2011production-note}（这也是芭蕾舞剧女主角的名字），
还有焰和使魔从玻璃门橱柜中破出的一幕，以及反复突显的发条机械，都使人联想到《胡桃夹子与老鼠王》：
胡桃夹子率领着从橱柜中走出的玩具军团，迎战钟声敲响时召出的老鼠王大军。

从最表层看，失去下颌的胡桃夹子意味着无用：它成了失却功能的物件。
然而，只要回到霍夫曼故事里嵌套叙事所讲述的胡桃夹子起源，这一意象便会显出更深层的含义。
胡桃夹子曾是被选中的人，一段类似预言称只有他能拯救遭鼠后诅咒的公主。
他必须完成一套复杂仪式才能救出公主，却偏偏在完成的瞬间被鼠后绊倒，诅咒反而落在自己身上。
这正像焰：她以为自己已经帮助小圆逃离命运，稍稍松懈，却最终发现丘比的介入令一切功亏一篑。
换言之，这又是她自责、自我惩罚的一种方式。

但是，魔法少女们拒绝配合。
她们不肯加入焰对焰的审判，不肯恨她，也不肯杀她。
相反，她们试图救出焰，打破迷宫和孵化者的圈套，以便小圆能够带她前往魔法少女的天国。
无论焰怎样疯狂呓语、苦苦哀求，她们仍执意宽恕她。
她们拒绝焰的判决，也要求焰拒绝它；
她们希望焰能原谅自己，并让自己获得自由。

但从影片开头起，焰便一直在与自己交战。
影片第一篇中，她追寻一个神秘而无形的暴君：
这个暴君统治着魔法少女们所在的、看似幸福的世界，而焰的意图是毁掉它。
她正是因为发现自己就是那名暴君，才对自己降下诅咒，彻底化为魔女；
这一切都是她对自身的叛逆。

这场叛逆直到影片结束也没有停止。
焰自称邪恶，并接受了衣着暴露、黑翼覆身的女恶魔形象。
但“我是邪恶的”与“我应当受罚”究竟有何区别？
这仍只是罪疚感的另一种表述，是她折磨自己的新方式。

她把自己抬升为宇宙级的存在——
一个看似拥有近乎无限力量、能够支配物质现实及其中之人的造物主：
她能改写沙耶香的记忆，令死者复生，又为小圆一家重构一段全新的历史，以逆转第一集的事件。
然而，她却选择创造一个令自己孑然一身的世界，与那些才刚开始和自己建立友谊的魔法少女隔绝开来。
她甚至选择在抹去沙耶香记忆前，先让沙耶香痛斥自己。

在焰创造的新现实中，她唯一真正流露出的情绪，是恐慌——
那发生在小圆几乎要重新与圆环之理相连之时。
换言之，那是小圆险些使一个象征希望与宽恕、能够终结焰痛苦的宇宙级存在重临之时。
焰尤其不能容许这种事发生。
她必须为没能救下小圆，又让小圆的境遇变得更糟而继续受苦。
她必须把小圆永远置于纯真与安全的状态，使她与自身的潜能隔绝。
因为“保护小圆”是焰唯一理解的善，而未能做到这一点，也就成为她唯一理解的恶。

原本一切都可以结束。
倘若其他魔法少女干脆杀了她，她便不会再受惩罚，痛苦也会随之终止。
然而，她们以残酷的仁慈强迫她活下去，迫使她另寻办法，一面继续保护小圆，一面继续惩罚自己。
她因此恨她们，恨她们不能像她恨自己那样恨她。
在新世界里，她以消极而带刺的举动嘲弄这些她憎恨的对象：
在麻美背后打碎茶杯，令人想起麻美与夏洛特交战时的死亡；
趁沙耶香的记忆逐渐消逝而嘲弄她，重演沙耶香化为魔女时的自我丧失；
又诱使杏子浪费食物。% TODOfig

而在片尾彩蛋中，她从白椅旁的悬崖纵身跃下，呼应早先小圆从椅子上翻落的画面。
她对自己的恨从未改变。
唯一改变的是，她如今有能力让魔法少女们也恨自己，将自己置于她们的对立面，期待她们替她完成那件事。

影片上映以来，围绕焰新身份的争论从未停歇：
她究竟是英雄，还是反派？
对此，答案是：对，焰既是《叛逆》的反派，也是与这个反派战斗的英雄。

但这个问题也可以有另一种答案：
不，焰其实是反派的受害者，而英雄必须将她救出。

她的魔女结界已经扩张至包覆整个宇宙。
如今，她就是整个故事。
